\def\stat{kruzhkov}





\def\tit{КОНЦЕПЦИЯ ПОСТРОЕНИЯ НАДКОРПУСНЫХ~БАЗ~ДАННЫХ$^*$}





\def\titkol{Концепция построения надкорпусных баз данных}





\def\aut{М.\,Г.~Кружков$^1$}





\def\autkol{М.\,Г.~Кружков}





\titel{\tit}{\aut}{\autkol}{\titkol}





\index{Кружков М.\,Г.}


\index{Kruzhkov M.\,G.}





{\renewcommand{\thefootnote}{\fnsymbol{footnote}}


\footnotetext[1]


{Работа выполнена в~Институте проблем информатики Федерального исследовательского центра 


<<Информатика и~управление>> Российской академии наук с~использованием ЦКП <<Информатика>> ФИЦ 


ИУ РАН.}} 





\renewcommand{\thefootnote}{\arabic{footnote}}


\footnotetext[1]{Федеральный исследовательский центр <<Информатика и~управление>> Российской академии наук, 


\mbox{magnit75@yandex.ru}}





 \label{st\stat}





 \thispagestyle{headings}


 


 \vspace*{-12pt} 





  


    


    \Abst{Представлен обзор концепции, основных структурных со\-став\-ля\-ющих и~функций 


надкорпусных баз данных (НБД). Надкорпусные базы данных пред\-став\-ля\-ют собой новый вид структурированных 


информационных ресурсов. Они существенно расширяют возможности лингвистических 


текс\-то\-вых корпусов, в~осо\-бен\-ности параллельных. Цель статьи~--- познакомить читателей 


с~основными возможностями параллельных корпусов, а~так\-же продемонстрировать, каким 


образом концепция НБД позволяет расширять эти возможности и~преодолевать некоторые 


их ограничения. В~рамках НБД лингвисты могут устанавливать, фиксировать 


и~аннотировать переводные соответствия (ПС) между языковыми единицами (ЯЕ) в~языках оригинала 


и~перевода, при этом для их аннотации используются руб\-ри\-ки фасетных классификаций, 


которые исследователи формируют в~соответствии со своими по\-треб\-но\-стя\-ми. 


Так\-же описана об\-щая архитектура НБД, разработанных в~ФИЦ ИУ РАН, которая 


под\-раз\-де\-ля\-ет\-ся на корпусную и~надкорпусную со\-став\-ля\-ющие, взаимодействующие друг 


с~другом в~рамках единой базы данных.}


    


    \KW{корпусная лингвистика; надкорпусная база данных; параллельный корпус; 


лингвистическое аннотирование; информационные технологии; фасетная классификация}





\DOI{10.14357\!/\!08696527210309} 





\vskip 10pt plus 9pt minus 6pt





 \thispagestyle{headings}


 


 \vspace*{-12pt}


    


\section{Введение}





\vspace*{-3pt}





   Корпуса параллельных текстов (параллельные корпуса) широко 


используются при проведении контрастивных (сопоставительных) 


лингвистических исследований, т.\,е.\ исследований, ставящих своей целью 


сравнение языковых явлений в~различных языках. Параллельные корпуса 


поз\-во\-ля\-ют давать ответы на различные вопросы сопоставительной 


лингвистики, основываясь на эмпирических данных большого объема, 


благодаря чему результаты контрастивных исследований во многом 


освобождаются от субъективности, которая может иметь мес\-то, когда 


исследователи обосновывают свои выводы, основываясь только на своей 


языковой интуиции. 


   


   Параллельные корпуса представляют собой собрания параллельных текс\-тов, 


каждый из которых вклю\-ча\-ет в~себя некоторый текст на языке оригинала (ЯО) 


и~один или несколько переводов этого текс\-та на другой язык (язык перевода, 


ЯП) или на несколько разных языков\footnote{В данной статье в~дальнейшем будет 


предполагаться, что параллельные тексты включают в~себя только один перевод~--- это позволит 


избежать постоянных оговорок (\textit{перевод} или \textit{переводы}, \textit{язык} или \textit{языки} 


и~т.\,д.). Особенности, возникающие при работе с~параллельными текстами, включающими в~себя 


более одного варианта перевода, будут оговариваться особо.}. В~параллельном тексте 


оригинал и~перевод разбиты на равное чис\-ло фраг\-мен\-тов, при этом каждому 


фраг\-мен\-ту в~текс\-те оригинала противопоставляется со\-от\-вет\-ст\-ву\-ющий ему по 


смыс\-лу фрагмент в~текс\-те перевода, благодаря чему в~структуре параллельного 


текс\-та образуются пары <<ори\-ги\-нал--пе\-ре\-вод>>. Про\-тя\-жен\-ность этих 


фраг\-мен\-тов в~раз\-ных корпусах, как правило, составляет от одного предложения 


до одного абзаца\footnote{Существуют также пословно выровненные тексты, однако они не 


представляют собой последовательность выровненных пар фрагментов, а~имеют более слож\-ную 


структуру, что вызвано тем, что разные языки имеют разный порядок слов, разные наборы 


служебных слов и~т.\,п.\ (см., например,~[1]).}. Для удобства восприятия очень длинные 


предложения иногда дробятся на несколько законченных по смыс\-лу 


фрагментов или, наоборот, несколько кратких предложений могут 


объединяться в~один расширенный фрагмент. Следует также учитывать, что 


границы предложений в~текс\-те оригинала и~текс\-те перевода могут не 


совпадать: одно предложение в~оригинале может быть разбито на два или 


больше предложений в~переводе или наоборот~--- одному предложению 


в~переводе может соответствовать два или больше предложений в~оригинале. 


   


   Кроме того, словоформы обычно снабжаются разметкой, помогающей 


лингвистам находить в~параллельных текс\-тах ин\-те\-ре\-су\-ющие их 


ЯЕ и~конструкции по морфологическим и\!/\!или синтаксическим 


приз\-на\-кам, а~так\-же посредством задания леммы~--- базовой формы 


ис\-сле\-ду\-емой единицы. По\-дроб\-нее о~структуре параллельных корпусов см., 


например,~[2].


   


   Корпуса параллельных текстов используются в~контрастивных 


лингвистических исследованиях, в~исследованиях по тео\-рии перевода, 


в~переводческой прак\-ти\-ке, а~так\-же в~учеб\-ных целях (см., например,~[3--5]). 


Однако работа с~параллельными корпусами за\-час\-тую требует значительных 


усилий от исследователей и~не всегда дает воз\-мож\-ность с~по\-мощью поисковых 


запросов напрямую получать ответы на поставленные вопросы, поскольку, 


составляя запросы, пользователь, как правило, не может достоверно 


установить, какие именно ЯЕ в~ЯП соответствуют искомым 


ЯЕ в~ЯО, и~наоборот. В~связи с~этим исследователям 


приходится анализировать значительное чис\-ло примеров вруч\-ную и~при этом 


либо удерживать в~памяти результаты этого анализа, либо фиксировать их на 


тех или иных носителях (в~текс\-то\-вых файлах, на бумажных карточках и~т.\,п.), 


чтобы впоследствии иметь воз\-мож\-ность сделать со\-от\-вет\-ст\-ву\-ющие обобщения 


на основе проанализированных данных. Однако зафиксированные по\-доб\-ным 


образом результаты анализа час\-то недостаточно хорошо структурированы; 


кроме того, в~дальнейшем их технически слож\-но обрабатывать, особенно когда 


речь идет о~значительных объемах языкового материала. Наконец, 


исследователи за\-час\-тую фиксируют таким образом не все ис\-сле\-ду\-емые 


единицы, а~только наиболее интересные и~показательные с~их точ\-ки зрения 


примеры, что не поз\-во\-ля\-ет сделать обобщенные выводы о~тех или иных 


явлениях, наблюдаемых во всем пред\-ста\-ви\-тель\-ном массиве параллельных 


текстов. Следовательно, такая организация работы не всегда позволяет 


в~пол\-ной мере использовать главные преимущества корпусного подхода: 


исследователи не име\-ют воз\-мож\-ности задействовать эффект мас\-шта\-ба, чтобы 


получить точ\-ные и~объективные количественные данные по способам 


и~особенностям перевода ис\-сле\-ду\-емых~ЯЕ. 


   


   Для решения этой проблемы в~Институте проблем информатики ФИЦ ИУ 


РАН была разработана и~реализована концепция <<надкорпусных баз 


данных>>. Надкорпусные базы данных призваны помочь исследователям выполнять 


мас\-штаб\-ные проекты в~об\-ласти сопоставительной лингвистики, опи\-ра\-ющи\-еся 


на богатый эмпирический материал, за\-клю\-чен\-ный в~параллельных  


корпусах~\cite{6-k, 7-k, 8-k, 9-k, 10-k, 11-k, 12-k, 13-k, 14-k, 15-k, 16-k}. Надкорпусные базы данных 


поз\-во\-ля\-ют группам исследователей анализировать большое чис\-ло примеров 


вместе с~их переводами и~\textit{сохранять результаты этого анализа 


в~унифицированном формате в~рам\-ках специально разработанной базы 


данных}, что позволяет в~дальнейшем лег\-ко анализировать полученный массив 


данных на количественном и~качественном уровне.


   


   Плунгян В.\,А.\ отмечает, что <<корпус$\ldots$ необходим 


профессиональным лингвистам, тем, кто так или иначе имеет дело с~фактами 


языка, а~значит, должен эти факты \textit{собирать 


и~систематизировать}>>~\cite{17-k} (курсив наш.~--- М.\,К.). Однако 


стандартный поисковый инструментарий большинства корпусов в~основном 


поз\-во\-ля\-ет лишь с~той или иной степенью точности находить в~корпусе 


определенные фак\-ты языка, но очень немногие корпуса предлагают функции, 


поз\-во\-ля\-ющие именно \textit{собирать и~систематизировать} эти факты, т.\,е.\ 


фиксировать их непосредственно в~структуре корпуса и~аннотировать их  


с~по\-мощью специальных классификационных признаков, с~тем чтобы 


впослед\-ст\-вии иметь воз\-мож\-ность получать обобщенную информацию по 


накопленному таким образом материалу. Именно эти функ\-ции стремятся взять 


на себя НБД, при этом очевидно, что они могли бы оказаться полезными и~при 


работе с~моноязычными корпусами. Однако именно при анализе параллельных 


корпусов концепция НБД оказывается в~осо\-бен\-ности полезной, поскольку она 


позволяет не только выделять и~аннотировать отдельные факты языка, но 


и~дает воз\-мож\-ность уста\-нав\-ли\-вать переводные соответствия между 


ЯЕ и~конструкциями в~определенных языковых парах, открывая новые 


возможности для исследователей в~об\-ласти контрастивной линг\-ви\-стики. 





%\vspace*{-12pt}





\section{Структура надкорпусной базы данных}





   В идеале структура НБД должна быть построена таким образом, чтобы 


пользователи имели возможность работать с~параллельным корпусом, находя 


в~нем употребления исследуемых ЯЕ и~их переводы, 


и~одновременно анализировать и~размечать выявленные таким образом 


переводные соответствия в~рамках заданной классификационной схемы. Этого 


мож\-но добиваться разными способами, однако общая архитектура сис\-те\-мы 


должна одновременно обеспечивать воз\-мож\-ность работы с~параллельным 


корпусом и~воз\-мож\-ность со\-хра\-не\-ния и~редактирования фор\-ми\-ру\-емых 


аннотаций переводных соответствий в~рамках базы данных. Информацию, 


которая генерируется на основе корпусных данных в~процессе аннотирования 


выявленных в~параллельном корпусе переводных соответствий, будем называть 


надкорпусной, чтобы отличать ее от собственно корпусной информации, 


за\-клю\-чен\-ной в~исходном параллельном корпусе.


   


   В целях обеспечения более удобного, быстрого и~надежного взаимодействия 


между корпусной и~надкорпусной составляющими было принято решение 


объединить их в~рамках единой базы данных. Далее будет описана структура 


НБД, которая разработана в~ФИЦ ИУ РАН и~объединяет в~себе корпусную 


и~надкорпусную со\-став\-ля\-ющие, а~так\-же обеспечивает механизмы 


взаимодействия между этими со\-став\-ля\-ющи\-ми. Следует принимать во 


внимание, что данная структура использовалась для реализации нескольких 


НБД, поддерживающих проекты различных кросслингвистических 


исследований~\cite{6-k, 7-k, 8-k, 9-k, 10-k, 11-k, 12-k, 13-k, 14-k, 15-k, 16-k}, 


вследствие чего некоторые детали этой общей схемы в~различных НБД могут 


незначительно отличаться.


   


\subsection{Корпусная составляющая}





   Корпусная составляющая предназначена для хранения параллельных 


текс\-тов, загруженных в~НБД. Источником текс\-тов служат параллельные 


корпуса Национального корпуса русского языка~[18]. В~этой час\-ти 


НБД содержится информация о~текс\-тах (произведениях), содержащихся 


в~корпусе, переводах этих текс\-тов, авторах текс\-тов и~их переводов, 


выровненных парах фрагментов в~рамках параллельных текс\-тов, 


словоупотреблениях, со\-став\-ля\-ющих эти фрагменты, и~знаках пунктуации. 


Также корпусная со\-став\-ля\-ющая содержит информацию о~возможных вариантах 


морфологического разбора всех включенных в~корпус словоформ, вклю\-чая 


указание их леммы (базовой формы слова) и~набора морфологических 


признаков. На основе этой информации пользователи НБД осуществляют 


первичный поиск фраг\-мен\-тов, в~которых могут находиться ин\-те\-ре\-су\-ющие их 


ЯЕ, после чего они получают воз\-мож\-ность формировать 


аннотации вы\-яв\-лен\-ных переводных соответствий. Логическая структурная 


схема корпусной со\-став\-ля\-ющей НБД пред\-став\-ле\-на на рис.~1.


   











  \smallskip


  


  %\noindent





\hangindent=5mm\hangafter=0\noindent{\small \textbf{Примечание.} Прямоугольники соответствуют таб\-ли\-цам НБД, линии между 


прямоугольниками~--- связям между таб\-ли\-ца\-ми. Связи <<один к~одному>> обозначаются 


прос\-ты\-ми линиями, связи <<один ко многим>>~--- линиями с~одним прос\-тым и~одним 


разветвленным концом, связи <<многие ко многим>>, ре\-а\-ли\-зу\-емые с~по\-мощью 


дополнительных промежуточных таб\-лиц,~--- линиями с~двумя раз\-ветв\-лен\-ны\-ми концами. 


Например, одной словоформе могут соответствовать несколько морфологических разборов, 


которые могут относиться к~одной или раз\-ным леммам (так, словоформе \textit{постели} 


соответствуют леммы \textit{постель} [сущ., род.\ падеж] и~\textit{постелить} [глагол, 


повелит.\ накл.]).}





\vspace*{-3pt}





\subsection{Надкорпусная составляющая}





\vspace*{-3pt}





\begin{figure} %fig1


\vspace*{1pt}


\begin{center}


 \mbox{%


 \epsfxsize=128mm 


\epsfbox{kru-1.eps}


 }


 \end{center}


\vspace*{-9pt}


\Caption{Логическая струк\-тур\-ная схема кор\-пус\-ной со\-став\-ля\-ющей НБД}


\vspace*{-6pt}


\end{figure}





   Структурная схема надкорпусной со\-став\-ля\-ющей НБД пред\-наз\-на\-че\-на для 


хранения аннотаций, выявленных в~параллельном корпусе ЯЕ,


и~ПС (рис.\,2).\linebreak Цент\-раль\-ный объ\-ект этой 


со\-став\-ля\-ющей~--- переводное


 соответствие\footnote{Во многих существующих 


работах вместо термина <<переводное соответствие>> используется термин моноэквиваленция (МЭ). 


Термин МЭ несет в~себе большую семантическую нагрузку, поскольку само его содержание требует 


определения того, что следует понимать под эк\-ви\-ва\-лент\-ностью применительно 


к~переводу. При этом в~любом случае приходится констатировать, что за\-час\-тую никакой <<эк\-ви\-ва\-лент\-ности>> 


в~зафиксированных ПС может быть не выявлено. Тем не менее термин МЭ 


традиционно широко используется применительно к~проб\-ле\-ма\-ти\-ке НБД, поэтому оба термина, 


ПС и~МЭ, в~на\-сто\-ящее время используются как синонимы.}. 


\begin{figure} %fig2


\vspace*{1pt}


\begin{center}


 \mbox{%


 \epsfxsize=128.084mm 


\epsfbox{kru-2.eps}


 }


 \end{center}


\vspace*{-9pt}


\Caption{Логическая структурная схема надкорпусной составляющей НБД}


\end{figure} 


Фраг\-мен\-ты ЯО и~ЯП, 


объединенные в~рамках ПС, описываются с~по\-мощью 


таб\-ли\-цы <<Фрагмент>>, при этом фраг\-мент ЯО интерпретируется 


и~фиксируется в~НБД как употреб\-ле\-ние некоторой языковой единицы ЯО, 


а~в~со\-от\-вет\-ст\-ву\-ющем ему фрагменте ЯП выявляется некоторая ЯЕ


или конструкция, которая отвечает за передачу исходной ЯЕ.


Примечательно, что один и~тот же фраг\-мент ЯО может быть задействован 


более чем в~одном ПС, если в~параллельном корпусе 


пред\-став\-ле\-но несколько раз\-ных переводов одного и~того же исходного \mbox{текста}.


{\looseness=1





}


   








   Языковые единицы ЯО и~ЯП фиксируются в~качестве классификационных 


рубрик фрагментов в~таб\-ли\-це <<Основные руб\-ри\-ки>>, при этом для удоб\-ст\-ва 


пользователя эти руб\-ри\-ки объединяются в~клас\-те\-ры. Кроме того, при анализе 


фрагментов исследователи присваивают фрагментам ЯО и~ЯП дополнительные 


классификационные руб\-ри\-ки, спе\-ци\-фи\-ци\-ру\-ющие различные аспекты 


упо\-треб\-ле\-ния рас\-смат\-ри\-ва\-емых ЯЕ, которые так\-же 


объединяются в~кластеры. Таким образом, фраг\-мен\-ты ЯО и~ЯП аннотируются 


с~по\-мощью \textit{фасетной классификации}. При этом следует отметить, что 


клас\-те\-ры основных и~дополнительных руб\-рик не всегда непосредственно 


соответствуют фактическим фасетам, что связано с~соображениями удоб\-ст\-ва 


пользователей НБД. Иногда несколько функционально или семантически 


близ\-ких фасетов группируются в~рам\-ках одного клас\-те\-ра: например, в~рам\-ках 


клас\-те\-ра <<Грамматические признаки>> объединяются руб\-ри\-ки нескольких 


фасетов, таких как \textit{время} (\textit{на\-сто\-ящее}, \textit{прошедшее}, 


\textit{будущее}), \textit{наклонение} (\textit{инфинитив}, \textit{императив}, 


\textit{причастие}, \textit{деепричастие}) и~т.\,д.


   


   Наконец, самим ПС так\-же присваиваются руб\-ри\-ки, 


специфицирующие те особенности ПС, которые 


определяют специфику переводной транс\-фор\-ма\-ции (например, 


\textit{перефразирование}, \textit{смена подлежащего} и~т.\,д.).





\section{Заключение}





   Как было отмечено выше, структура НБД должна обеспечивать воз\-мож\-ность 


   параллельной работы и~с корпусной, и~с~надкорпусной со\-став\-ля\-ющи\-ми, 


в~связи с~чем было принято решение объ\-еди\-нить их в~рамках одной базы 


данных. Как видно из рис.~1 и~2, эти составляющие связывает таб\-ли\-ца 


<<Сло\-во\-упо\-треб\-ле\-ние>>. Между таб\-ли\-ца\-ми <<Фрагмент>> 


и~<<Сло\-во\-упо\-треб\-ле\-ние>> находится свя\-зы\-ва\-ющая таб\-ли\-ца, которая явно 


специфицирует, какие сло\-во\-упо\-треб\-ле\-ния из параллельного корпуса входят во 


фраг\-мен\-ты ЯО и~ЯП, служащие формообразующими элементами 


ПС. Данная свя\-зы\-ва\-ющая таб\-ли\-ца так\-же позволяет 


выделять в~рамках этих фрагментов сами рас\-смат\-ри\-ва\-емые ЯЕ 


и~<<функционально значимые слова>>, т.\,е.\ слова, непосредственно 


влия\-ющие на функционирование этих ЯЕ. По\-дроб\-нее процесс 


по\-стро\-ения ПС на основе кор\-пус\-ной информации описан в~работе~[19].


   


   Возможность обращения к~корпусной информации так\-же играет важ\-ную 


роль при работе с~массивом уже по\-стро\-ен\-ных аннотаций ПС. В~част\-ности, 


корпусная информация используется в~сис\-те\-ме просмотра и~поиска 


ПС, поз\-во\-ляя осуществлять поиск аннотаций по\-сред\-ст\-вом задания 


автора исходного произведения или базовой формы тех или иных лексем 


(по\-дроб\-нее о~поисковых функциях НБД см., например,~[20]). Наконец, 


корпусная информация так\-же может быть задействована при по\-стро\-ении 


запросов к~массиву сформированных ПС (см., 


например,~[21]). 


   


{\small\frenchspacing


{%\baselineskip=10.8pt


\addcontentsline{toc}{section}{Литература}


\begin{thebibliography}{99}





\bibitem{1-k}


\Au{Морозова Ю.\,И., Козеренко Е.\,Б., Шарнин~М.\,М.} Методика извлечения по\-слов\-ных 


переводных соответствий из параллельных текс\-тов с~применением моделей дистрибутивной 


семантики~/\!\!/ Сис\-те\-мы и~средства информатики, 2014. Т.~24. №\,2. С.~131--142.


\bibitem{2-k}


\Au{Кружков М.\,Г.} Информационные ресурсы контрастивных лингвистических 


исследований: электронные корпуса текстов~/\!\!/ Сис\-те\-мы и~средства информатики, 2015. 


Т.~25. №\,2. С.~140--159.





\bibitem{5-k} %3


\Au{Johansson S.} Seeing through Multilingual Corpora: On the use of corpora in contrastive 


studies.~--- Studies in corpus linguistics ser.~--- Amsterdam/Philadelphia: John Benjamins Publishing Co., 2007. Vol.~26. 378~p.





\bibitem{3-k} %4


\Au{Добровольский Д.\,О., Кротова~Е.\,Б., Парина~И.\,С.} Корпусная лексикография 


(материалы мас\-тер-клас\-са)~/\!\!/ Русская германистика: ежегодник Российского союза 


германистов.~--- М.: Языки славянской культуры,  2014. Т.~11. С.~237--278.


\bibitem{4-k} %5


\Au{Сичинава Д.\,В.} Параллельные корпуса в~со\-ста\-ве Национального корпуса русского 


языка: новые языки и~новые задачи~/\!\!/ Труды Института русского языка им.\ 


В.\,В.~Виноградова, 2019. Вып.~21. С.~41--60.








\bibitem{6-k}


\Au{Зализняк Анна~А., Зацман~И.\,М., Инькова~О.\,Ю., Кружков~М.\,Г.} Надкорпусные базы 


данных как лингвистический ресурс~/\!\!/ Корпусная лингвистика: Тр. VII Междунар. 


конф.~--- СПб.: СПбГУ, 2015. С.~211--218.


\bibitem{7-k}


Семантика коннекторов: контрастивное исследование~/ Под. ред. О.\,Ю.~Иньковой.~--- М.: 


ТОРУС ПРЕСС, 2018. 368~с.


\bibitem{8-k}


\Au{Зацман И.\,М., Кружков~М.\,Г.} Надкорпусная база данных коннекторов: развитие 


сис\-те\-мы терминов проектирования~/\!\!/ Сис\-те\-мы и~средства информатики, 2018. Т.~28. 


№\,4. С.~156--167.


\bibitem{9-k}


\Au{Бунтман Н.\,В., Инькова О.\,Ю.} Французский коннектор \textit{sinon}~/\!\!/ Семантика 


коннекторов: контрастивное исследование~/ Под ред.\ О.\,Ю.~Иньковой.~--- М.: ТОРУС 


ПРЕСС, 2018. С.~301--330.





\bibitem{16-k} %10


\Au{Добровольский Д.\,О., Зализняк Анна~А.} Немецкие конструкции с~модальными 


глаголами и~их рус\-ские соответствия: проект надкорпусной базы данных~/\!\!/ 


Компьютерная лингвистика и~интеллектуальные технологии: По мат-лам Междунар. конф. 


<<Диалог>>.~--- М.: РГГУ, 2018. С.~172--184.











\bibitem{11-k}


Структура коннекторов и~методы ее описания~/ Под. ред. О.\,Ю.~Иньковой.~--- М.: ТОРУС 


ПРЕСС, 2019. 316~с.


\bibitem{12-k}


\Au{Зацман~И.\,М., Кружков М.\,Г., Лощилова~Е.\,Ю.} Методы и~средства информатики для 


описания струк\-ту\-ры неоднословных коннекторов~/\!\!/ Структура коннекторов и~методы ее 


описания~/ Под ред. О.\,Ю.~Иньковой.~--- М.: ТОРУС ПРЕСС, 2019. С.~205--230.


\bibitem{13-k}


\Au{Гончаров А.\,А., Инькова О.\,Ю., Кружков~М.\,Г.} Методология аннотирования 


в~надкорпусных базах данных~/\!\!/ Сис\-те\-мы и~средства информатики, 2019. Т.~29. №\,2. 


С.~148--160.





\bibitem{10-k} %14


\Au{Инькова О.\,Ю., Кружков~М.\,Г.} Со\-че\-та\-емость ло\-ги\-ко-се\-ман\-ти\-че\-ских отношений: 


количественные методы анализа~/\!\!/ Информатика и~её применения, 2019. Т.~13. Вып.~2. 


С.~83--91.





\bibitem{14-k} %15


Семантика коннекторов: количественные методы описания~/ Под ред. О.~Иньковой.~--- 


Bern\!/\!Berlin: Peter Lang AG, 2021. 276~с.


\bibitem{15-k} %16


\Au{Гончаров А.\,А., Бунтман~Н.\,В., Нуриев~В.\,А.} Ошибки в~машинном переводе 


коннекторов: сравнительный анализ работы двух автоматических переводчиков~/\!\!/ 


Семантика коннекторов. Количественные методы описания~/ Под ред. О.~Иньковой.~--- 


Bern\!/\!Berlin: Peter Lang AG, 2021. С.~225--276.


\bibitem{17-k}


\Au{Плунгян В.\,А.} Зачем нужен Национальный корпус русского языка? Неформальное 


введение~/\!\!/ Национальный корпус русского языка: 2003--2005.~--- М.: Индрик, 2005.  


С.~6--20.


\bibitem{18-k}


Национальный корпус русского языка. {\sf https://ruscorpora.ru/new/index.html.}


\bibitem{19-k}


\Au{Попкова Н.\,А., Инькова О.\,Ю., Зацман~И.\,М., Кружков~М.\,Г.} Методика по\-стро\-ения 


моноэквиваленций в~надкорпусной базе данных коннекторов~/\!\!/ Задачи современной 


информатики: Тр.\ 2-й Молодежной научн. конф.~--- М.: ФИЦ ИУ РАН, 2015. С.~143--153.


\bibitem{20-k}


\Au{Инькова О.\,Ю., Кружков~М.\,Г.} Надкорпусные рус\-ско-фран\-цуз\-ские базы данных 


глагольных форм и~коннекторов~/\!\!/ Langues slaves en contraste.~--- Eds.\ O.~Inkova, 


A.~Trovesi.~--- Bergamo: Bergamo University Press, 2016. P.~365--392.


\bibitem{21-k}


\Au{Зацман И.\,М., Кружков М.\,Г., Лощилова~Е.\,Ю.} Методы анализа час\-тот\-ности моделей 


перевода коннекторов и~об\-ра\-ти\-мость генерализации статистических моделей~/\!\!/ Системы 


и~средства информатики, 2017. Т.~27. №\,4. С.~164--176.


  \end{thebibliography}


} }











\vspace*{-4pt}





\hfill{\small\textit{Поступила в~редакцию 14.08.21}}





\vspace*{6pt}





\hrule





\vspace*{2pt}





%\newpage





%\vspace*{-28pt}





%\hrule





%\vspace*{8pt}





\def\tit{CONCEPTUAL FRAMEWORK FOR~SUPRACORPORA DATABASES}





\def\titkol{Conceptual framework for~supracorpora databases}





\def\aut{M.\,G.~Kruzhkov}





\def\autkol{M.\,G.~Kruzhkov}





\titel{\tit}{\aut}{\autkol}{\titkol}





\vspace*{-11pt}





\noindent


Federal Research Center ``Computer Science and Control'' of the Russian


 Academy of Sciences, 44-2~Vavilov Str., Moscow 119133, Russian Federation





\def\leftfootline{\small{\textbf{\thepage}


\hfill Sistemy i Sredstva Informatiki~---


Systems and Means of Informatics 2021 vol~31 no\,3}


}%


 \def\rightfootline{\small{Sistemy i Sredstva Informatiki~---


 Systems and Means of Informatics 2021 vol~31 no\,3


\hfill \textbf{\thepage}}}





\vspace*{3pt}











\Abste{The paper provides an overview of the concept, main structural 


constituents, and functions of supracorpora databases (SCDB). 


Supracorpora databases represent a~novel type of structured information resources that 


significantly expand capabilities of linguistic text corpora, parallel 


corpora in particular. The paper outlines principle features and limitations of parallel 


corpora and demonstrates how SCDBs allow extending these features and overcoming 


the limitations. Supracorpora databases allow linguistic experts to establish, record, and annotate 


translation correspondences between language units in the source and target texts 


while relying on faceted classification categories composed by the researchers 


themselves according to their requirements. The article also describes the general 


structure of SCDB architecture developed in FRC CSC RAS which incorporates corpus and 


subcorpus constituents that interact with one another as a~part of a~common database.}





\KWE{corpus linguistics; supracorpora database; parallel corpus; linguistic annotation; information 


technologies; faceted classification}











\DOI{10.14357\!/\!08696527210309}  





%\vspace*{-24pt}





\Ack


\noindent


The research was carried out using the infrastructure of the Shared Research Facilities ``High Performance 


Computing and Big Data'' (CKP ``Informatics'') of FRC CSC RAS (Moscow).








%\vspace*{-12pt}





\renewcommand{\bibname}{\protect\rmfamily References}


%\renewcommand{\bibname}{\large\protect\rm References}





{\small\frenchspacing


{%\baselineskip=10.8pt


\addcontentsline{toc}{section}{References}


\begin{thebibliography}{99}


\bibitem{1-k-1}


\Aue{Morozova, Yu.\,I., E.\,B.~Kozerenko, and M.\,M.~Sharnin.} 2014.Metodika izvlecheniya poslovnykh 


perevodnykh sootvetstviy iz parallel'nykh tekstov s~primeneniem modeley distributivnoy semantiki [Method 


for extracting single-word translation corresponcences from parallel texts using distributional semantics 


models]. \textit{Sistemy i~Sredstva Informatiki~--- Systems and Means of Informatics} 24(2):131--142.


\bibitem{2-k-1}


\Aue{Kruzhkov, M.\,G.} 2015. Informatsionnye resursy kontrastivnykh lingvisticheskikh issledovaniy: 


elektronnye korpusa tekstov [Information resources for contrastive studies: Electronic text corpora]. 


\textit{Sistemy i~Sredstva Informatiki~--- Systems and Means of Informatics} 25(2):140--159.





\bibitem{5-k-1} %3


\Aue{Johansson, S.} 2007. \textit{Seeing through Multilingual Corpora: On the use of corpora in contrastive 


studies}. Studies in corpus linguistics ser.


Amsterdam\!/\!Philadelphia: John Benjamins Publishing Co. Vol.~26. 378~p. 





\bibitem{3-k-1} %4


\Aue{Dobrovol'skiy, D.\,O., E.\,B.~Krotova, and I.\,S.~Parina.} 2014. Korpusnaya leksikografiya (materialy 


master-klassa) [Corpus lexicography (workshop materials)]. \textit{Russkaya germanistika: ezhegodnik 


Rossiyskogo soyuza germanistov} [Russian German studies: Yearbook of the Russian Union of Germanists]. 


Moscow: Yazy\-ki sla\-vyan\-skoy kul'tury. 11:237--278.


\bibitem{4-k-1} %5


\Aue{Sichinava, D.\,V.} 2019. Parallel'nye korpusa v~so\-sta\-ve Natsional'nogo korpusa russkogo yazyka: 


novye yazyki i~novye zadachi [Parallel corpora within the Russian National Corpus: New languages and new 


objectives]. \textit{Trudy Instituta russkogo yazyka im.\ V.\,V.~Vinogradova} [V.\,V.~Vinogradov Russian 


Language Institute Proceedings]. 21:41--60.





\bibitem{6-k-1}


\Aue{Zaliznyak, A.\,A., I.\,M.~Zatsman, O.\,Yu.~Inkova, and M.\,G.~Kruzhkov.} 2015. Nadkorpusnye 


bazy dannykh kak lingvisticheskiy resurs [Supracorpora databases as linguistic resource]. \textit{Tr.\ 7-y 


Mezhdunar. konf.  ``Korpusnaya lingvistika''} [7th Conference (International) on Corpus Linguistics 


Proceedings]. St.\ Petersburg: St.\ Petersburg State University. 211--218.


\bibitem{7-k-1}


Inkova, O., ed. 2018. \textit{Semantika konnektorov: kontrastivnoe issledovanie} [Semantics of connectives: 


A~contrastive study]. Moscow: TORUS PRESS. 368~p.


\bibitem{8-k-1}


\Aue{Zatsman, I.\,M., and M.\,G.~Kruzhkov.} 2018. Nadkorpusnaya baza dannykh konnektorov: razvitie 


sistemy terminov proektirovaniya [Supracorpora database of connectives: Design-oriented evolution of the 


term system]. \textit{Sistemy i~Sredstva Informatiki~--- Systems and Means of Informatics} 28(4):156--167.


\bibitem{9-k-1}


\Aue{Buntman, N.\,V., and O.\,Yu.~Inkova}. 2018. Frantsuzskiy konnektor \textit{sinon} [French connective 


\textit{sinon}]. \textit{Semantika konnektorov: kontrastivnoe issledovanie} [Semantics of connectives: A~contrastive 


study]. Ed.\,O.\,Yu.~Inkova. Moscow: TORUS PRESS. 301--330.





\bibitem{16-k-1} %10


\Aue{Dobrovol'skiy, D.\,O., and Anna A.~Zaliznyak.} 2018. Nemetskie konstruktsii s~mo\-dal'\-ny\-mi 


glagolami i~ikh russkie sootvetstviya: proekt nadkorpusnoy bazy dannykh [German constructions with modal 


verbs and their Russian correlates: A~supracorpora database project]. \textit{Komp'yuternaya lingvistika 


i~intellektual'nye tekhnologii. Po mat-lam ezhegodnoy Mezhdunar. konf. ``Dialog''} [Computational 


Linguistics and Intellectual Technologies. Papers from the Annual Conference (International) ``Dialogue'']. 


Moscow: RGGU. 172--184.








\bibitem{11-k-1}


Inkova, O., ed. 2019. \textit{Struktura konnektorov i~metody ee opi\-sa\-niya} [Structure of connectives and 


methods for its description]. Moscow: TORUS PRESS. 316~p.


\bibitem{12-k-1}


\Aue{Zatsman, I., M.~Kruzhkov, and E.~Loshchilova.} 2019. Metody i~sredstva informatiki dlya opisaniya 


struktury neodnoslovnykh konnektorov [Methods and means of informatics for multiword connectives 


structure description]. \textit{Struktura konnektorov i~metody ee opisaniya} [Connectives structure and 


methods of its description]. Ed. O.\,Yu.~Inkova. Moscow: TORUS PRESS. 205--230.


\bibitem{13-k-1}


\Aue{Goncharov, A.\,A., O.\,Yu.~Inkova, and M.\,G.~Kruzhkov} 2019. Metodologiya annotirovaniya 


v~nadkorpusnykh bazakh dannykh [Annotation methodology of supracorpora databases]. \textit{Sistemy 


i~Sredstva Informatiki~--- Systems and Means of Informatics} 29(2):148--160.





\bibitem{10-k-1} %14


\Aue{Inkova, O.\,Yu., and M.\,G.~Kruzhkov.} 2019. Sochetaemost' logiko-semanticheskikh otnosheniy: 


kolichestvennye metody analiza [Compatibility of logical semantic relations: Methods of quantitative 


analysis]. \textit{Informatika i~ee Primeneniya~--- Inform. Appl}. 13(2):83--91.





\bibitem{14-k-1} %15


Inkova, O., ed. 2021. \textit{Semantika konnektorov: kolichestvennye metody opisaniya} [Semantics of 


connectives: Quantitative methods of analysis]. Bern\!/\!Berlin: Peter Lang. 276~p.


\bibitem{15-k-1} %16


\Aue{Goncharov, A.\,A., N.\,V.~Buntman, and V.\,A.~Nuriev.} 2021. Oshibki v~mashinnom perevode 


konnektorov: sravnitel'nyy analiz raboty dvukh avtomaticheskikh perevodchikov [Errors in machine 


translation of connectives: Comparative analysis of two translation engines]. \textit{Semantika konnektorov: 


kolichestvennye metody opisaniya} [Semantics of connectives: Quantitative methods of analysis].


Ed.\ O.~Inkova.  Bern\!/\!Berlin: Peter Lang AG. 225--276.








\bibitem{17-k-1}


\Aue{Plungyan, V.\,A.} 2005. Zachem nuzhen Natsional'nyy korpus russkogo yazyka? Neformal'noe 


vvedenie [What the Russian National Corpus is for? Informal introduction]. \textit{Natsional'nyy korpus 


russkogo yazyka: 2003--2005} [Russian National Corpus 2003--2005]. Moscow: Indrik. 6--20.


\bibitem{18-k-1}


Natsional'nyy korpus russkogo yazyka [Russian National Corpus]. Available at: {\sf 


http:// www.ruscorpora.ru} (accessed August~11, 2021).


\bibitem{19-k-1}


\Aue{Popkova, N.\,A., O.\,Yu.~Inkova, I.\,M.~Zatsman, and M.\,G.~Kruzhkov.} 2015. Metodika 


postroeniya monoekvivalentsiy v~nadkorpusnoy baze dannykh konnektorov [Methodology of constructing 


monoequivalences in the supracorpora database of connectors]. \textit{Tr. 2-y Molodezhnoy nauchn. konf. 


``Zadachi sovremennoy informatiki''} [2nd Scientific Conference ``Problems of Modern Informatics'' 


Proceedings]. Moscow: FRC CSC RAS. 143--153.


\bibitem{20-k-1}


\Aue{Inkova, O.\,Yu., and M.\,G.~Kruzhkov.} 2016. Nadkorpusnye russko-frantsuzskie basy dannykh 


glagol'nykh form konnektorov [Supracorpora databases of Russian and French verbal forms and connectors]. 


\textit{Lingue slave a~confronto}. Eds. O.~Inkova and A.~Trovesi. Bergamo: Bergamo University Press. 


365--392.


\bibitem{21-k-1}


\Aue{Zatsman, I.\,M., M.\,G. Kruzhkov, and E.\,Yu.~Loshchilova.} 2017. Metody analiza chas\-tot\-nosti 


modeley perevoda konnektorov i~ob\-ra\-ti\-most' ge\-ne\-ra\-li\-za\-tsii 


statisticheskikh dan\-nykh [Methods of frequency 


analysis of connectives translations and reversibility of statistical data generalization]. \textit{Sistemy 


i~Sredstva Informatiki~--- Systems and Means of Informatics} 27(4):164--176.


\end{thebibliography}


} }





%\vspace*{-3pt}





\hfill{\small\textit{Received August 14, 2021}} 





%\vspace*{-16pt} 





\Contrl





\noindent


\textbf{Kruzhkov Mikhail G.} (b.\ 1975)~--- senior scientist, Institute of Informatics Problems, Federal 


Research Center ``Computer Science and Control'' of the Russian Academy of Sciences, 44-2~Vavilov Str., 


Moscow 119333, Russian Federation; \mbox{magnit75@yandex.ru}








\label{end\stat}





\renewcommand{\bibname}{\protect\rm Литература} 


